% !TeX program = pdflatex
\documentclass[12pt,a4paper,oneside]{article}

% PAQUETES
\usepackage[utf8]{inputenc}
\usepackage[T1]{fontenc}
\usepackage[spanish,es-tabla]{babel}
\usepackage{newtxtext,newtxmath}
\usepackage{geometry}
\usepackage{array}
\usepackage{longtable}
\usepackage{booktabs}
\usepackage{enumitem}
\usepackage{microtype}
\usepackage{setspace}
\usepackage{hyperref}
\usepackage{titlesec}

\geometry{
    left=2.5cm,
    right=2.5cm,
    top=2.5cm,
    bottom=2.5cm
}

\setstretch{1.15}
\setlength{\parindent}{1.25cm}
\setlength{\parskip}{4pt}
\renewcommand{\arraystretch}{1.25}
\setlength{\tabcolsep}{3pt}

\hypersetup{
    colorlinks=true,
    linkcolor=black,
    citecolor=black,
    urlcolor=blue,
    pdftitle={FORM-TL-01 - Taller de tipos de requisitos},
    pdfauthor={Angelina Julexy Llerena Abril}
}

\titleformat{\section}
{\normalfont\bfseries\large}{\thesection.}{0.5em}{}

\titleformat{\subsection}
{\normalfont\bfseries\normalsize}{\thesubsection.}{0.5em}{}

% Columnas para las tablas
\newcolumntype{C}[1]{>{\centering\arraybackslash}p{#1}}
\newcolumntype{L}[1]{>{\raggedright\arraybackslash}p{#1}}

% DATOS
\newcommand{\nombreestudiante}{Angelina Julexy Llerena Abril}
\newcommand{\paralelo}{A}
\newcommand{\fechaentrega}{31 de octubre de 2026}

% DOCUMENTO
\begin{document}

% PORTADA
\begin{titlepage}
    \centering

    {\large \textbf{UNIVERSIDAD TÉCNICA ESTATAL DE QUEVEDO}}\\[0.3cm]
    {\large Facultad de Ciencias de la Computación}\\
    {\large Carrera de Software}\\[1.5cm]

    {\Large \textbf{INGENIERÍA DE REQUERIMIENTOS}}\\[0.5cm]
    {\large ISR-401}\\[1.5cm]

    {\LARGE \textbf{FORM-TL-01}}\\[0.4cm]
    {\Large Taller de tipos de requisitos}\\[0.5cm]

    \rule{0.85\textwidth}{0.5pt}\\[0.5cm]
    {\large \textbf{Identificación y clasificación de requisitos}}\\
    {\large Plataforma de gestión de citas médicas}\\[0.5cm]
    \rule{0.85\textwidth}{0.5pt}\\[2cm]

    \begin{flushleft}
        \textbf{Estudiante:} \nombreestudiante\\[0.4cm]
        \textbf{Nivel:} Cuarto\\[0.4cm]
        \textbf{Paralelo:} \paralelo\\[0.4cm]
        \textbf{Docente:} Gleiston Cicerón Guerrero Ulloa, PhD\\[0.4cm]
        \textbf{Fecha de entrega:} \fechaentrega
    \end{flushleft}

    \vfill
    {\large Quevedo -- Ecuador}\\
    {\large 2026}
\end{titlepage}

% CONTENIDO
\section{Clasificación de los requisitos}

La clasificación de los 20 enunciados considera
tres dimensiones: el tipo de requisito, su naturaleza y el nivel
al que corresponde. El tipo comprende las categorías
funcional, calidad y restricción, de acuerdo con lo establecido
en la guía entregada por el docente. La naturaleza distingue entre
requisitos del producto y del proceso, mientras que el nivel
permite diferenciar los correspondientes al sistema y al software.

Los criterios de clasificación se apoyan en
ISO/IEC/IEEE 29148:2018 \cite{ieee29148}. Para los requisitos de
calidad también se consideran las características y subcaracterísticas
definidas en ISO/IEC 25010:2023 \cite{iso25010}.

\subsection{Tabla de clasificación}

\small
\begin{longtable}{
    C{1.25cm}
    L{3.5cm}
    L{2.2cm}
    C{1.7cm}
    C{1.35cm}
    L{4.5cm}
}
\caption{Clasificación de los 20 enunciados del caso de estudio}
\label{tab:clasificacion}\\

\toprule
\textbf{Código} &
\textbf{Enunciado} &
\textbf{Tipo} &
\textbf{Naturaleza} &
\textbf{Nivel} &
\textbf{Justificación} \\
\midrule
\endfirsthead

\multicolumn{6}{c}{
    
}\\
\toprule
\textbf{Código} &
\textbf{Enunciado} &
\textbf{Tipo} &
\textbf{Naturaleza} &
\textbf{Nivel} &
\textbf{Justificación} \\
\midrule
\endhead

\bottomrule
\endfoot
E01 &
El paciente podrá reservar una cita eligiendo especialidad,
médico, fecha y hora entre los horarios disponibles. &
Funcional &
Producto &
Software &
Según ISO/IEC/IEEE 29148:2018, corresponde a un requisito
funcional porque especifica el servicio de reserva de citas
que debe proporcionar el software \cite{ieee29148}. \\

E02 &
La plataforma enviará al paciente un recordatorio por mensaje
de texto 24 horas antes de su cita. &
Funcional &
Producto &
Software &
Establece un servicio que debe proporcionar el software,
por lo que corresponde a un requisito funcional según
ISO/IEC/IEEE 29148:2018 \cite{ieee29148}. \\

E03 &
La aplicación debe ser rápida. &
Calidad: eficiencia del desempeño
(comportamiento temporal) &
Producto &
Software &
Corresponde a eficiencia del desempeño porque establece
una condición sobre el tiempo de respuesta del software,
característica definida en ISO/IEC 25010:2023
\cite{iso25010}. \\

E04 &
El recepcionista podrá registrar la llegada del paciente que
se acerca a la ventanilla. &
Funcional &
Producto &
Software &
Es un requisito funcional porque especifica la operación
de registro de llegada que debe ofrecer el software
\cite{ieee29148}. \\

E05 &
La plataforma se desarrollará en Java 21 con base de datos
PostgreSQL, tecnologías que domina el área de TI. &
Restricción &
Producto &
Software &
Impone tecnologías que condicionan la solución de software,
por lo que constituye una restricción según
ISO/IEC/IEEE 29148:2018 \cite{ieee29148}. \\

E06 &
El médico podrá consultar la lista de pacientes del día y
marcar cada consulta como atendida. &
Funcional &
Producto &
Software &
Define las funciones de consulta de pacientes y registro
de atenciones que debe proporcionar el software
\cite{ieee29148}. \\

E07 &
Los datos de salud de los pacientes deben estar seguros. &
Calidad: seguridad (confidencialidad) &
Producto &
Software &
Establece una condición de protección de la información
frente a accesos no autorizados, correspondiente a la
confidencialidad de ISO/IEC 25010:2023 \cite{iso25010}. \\

E08 &
La aplicación móvil debe funcionar en teléfonos con Android
10 o superior. &
Restricción &
Producto &
Software &
Fija una condición de compatibilidad que limita los entornos
de ejecución del software \cite{ieee29148}. \\

E09 &
Una cita cancelada con menos de cuatro horas de anticipación
se registrará como inasistencia. &
Funcional &
Producto &
Software &
Especifica la respuesta que debe ejecutar el software
ante una cancelación fuera del plazo establecido
\cite{ieee29148}. \\

E10 &
La plataforma debe ser fácil de usar para los pacientes
adultos mayores. &
Calidad: capacidad de interacción
(inclusividad) &
Producto &
Software &
Establece una condición de inclusión para distintos grupos
de usuarios, característica contemplada en
ISO/IEC 25010:2023 \cite{iso25010}. \\

E11 &
El proyecto debe entregarse en un plazo máximo de seis meses. &
Restricción &
Proceso &
Software &
Establece un límite temporal que condiciona el proceso
de desarrollo y entrega del software \cite{ieee29148}. \\

E12 &
La sesión de un usuario se cerrará automáticamente tras
10 minutos sin actividad. &
Funcional &
Producto &
Software &
Define la función de cierre automático de sesión que
debe ejecutar el software ante la inactividad
\cite{ieee29148}. \\

E13 &
El quiosco de la sala de espera leerá el número de cédula
ingresado en su pantalla táctil e imprimirá un ticket con
el número de turno. &
Funcional &
Producto &
Sistema &
Especifica una capacidad funcional del sistema que integra
la lectura de datos y la impresión de turnos
\cite{ieee29148}. \\

E14 &
La plataforma debe estar siempre disponible. &
Calidad: fiabilidad (disponibilidad) &
Producto &
Sistema &
Establece una condición de disponibilidad operativa
del sistema, subcaracterística de fiabilidad definida
en ISO/IEC 25010:2023 \cite{iso25010}. \\

E15 &
El director médico podrá generar un reporte mensual de
inasistencias por especialidad. &
Funcional &
Producto &
Software &
Define el servicio de generación de reportes que debe
proporcionar el software al director médico
\cite{ieee29148}. \\

E16 &
Cada incremento se validará con el director médico al final
de cada iteración de dos semanas. &
Restricción &
Proceso &
Software &
Establece una condición obligatoria para las actividades
de validación durante el proceso de desarrollo
\cite{ieee29148}. \\

E17 &
El tratamiento de los datos personales de los pacientes debe
cumplir la Ley Orgánica de Protección de Datos Personales
del Ecuador. &
Restricción &
Producto &
Sistema &
Impone una obligación legal que condiciona el funcionamiento
y tratamiento de datos del sistema \cite{ieee29148}. \\

E18 &
La plataforma debe integrarse sin problemas con el sistema
de facturación actual. &
Calidad: compatibilidad (interoperabilidad) &
Producto &
Software &
Establece el intercambio de información con otro producto,
propio de la interoperabilidad definida en
ISO/IEC 25010:2023 \cite{iso25010}. \\

E19 &
Si el servidor principal falla, la plataforma debe restablecer
el servicio en menos de 30 minutos sin perder ninguna cita
confirmada. &
Calidad: fiabilidad (recuperabilidad) &
Producto &
Sistema &
Define la recuperación del sistema ante fallos y la
conservación de los datos, características contempladas
en ISO/IEC 25010:2023 \cite{iso25010}. \\

E20 &
Todo cambio en el código fuente será revisado por otro
integrante del equipo antes de integrarse a la rama principal. &
Restricción &
Proceso &
Software &
Establece una práctica obligatoria de revisión dentro
del proceso de desarrollo del software \cite{ieee29148}. \\

\end{longtable}
\normalsize

\subsection{Criterios utilizados en la clasificación}

Las dimensiones de la tabla se interpretan mediante
los siguientes criterios, de acuerdo con la guía del docente
y las normas de referencia:

\begin{itemize}[leftmargin=1.5cm]

    \item \textbf{Funcional:} describe una función, acción o
    servicio que el sistema debe realizar para responder a
    las necesidades de sus usuarios.

    \item \textbf{Calidad:} establece una característica que
    determina cómo debe comportarse el producto o qué
    propiedades debe cumplir. Se identifica mediante las
    características y subcaracterísticas de ISO/IEC 25010:2023.

    \item \textbf{Restricción:} representa una condición,
    limitación o regla obligatoria que condiciona la solución
    o el proceso utilizado para desarrollarla.

    \item \textbf{Producto:} corresponde a los requisitos
    relacionados con las características, funciones o
    condiciones que debe cumplir el producto desarrollado.

    \item \textbf{Proceso:} comprende los requisitos que
    establecen condiciones sobre las actividades, prácticas
    o procedimientos utilizados durante el desarrollo.

    \item \textbf{Sistema:} se refiere a requisitos que abarcan
    el funcionamiento de la solución en su conjunto, incluidos
    los componentes de hardware, software y sus interacciones.

    \item \textbf{Software:} corresponde a requisitos relacionados
    específicamente con las funciones o características de la
    aplicación y sus componentes de software.

\end{itemize}

Las justificaciones de la Tabla~\ref{tab:clasificacion}
explican los motivos de cada clasificación. Se consideran
los criterios de ingeniería de requisitos de
ISO/IEC/IEEE 29148:2018 \cite{ieee29148}; para los requisitos
de calidad también se toma en cuenta la característica
y subcaracterística correspondiente de ISO/IEC 25010:2023
\cite{iso25010}.

% EVALUACIÓN DE PROPIEDADES DE CALIDAD
\section{Enunciados con propiedades de calidad violadas y su cuantificación}

En el caso de estudio se identifican cinco requisitos que
presentan problemas en sus propiedades de calidad,
principalmente ambigüedad y falta de verificabilidad:
E03, E07, E10, E14 y E18. Sus formulaciones contienen
expresiones generales que dificultan determinar con precisión
cuándo se considera cumplido cada requisito.

\begin{table}[htbp]
\centering
\caption{Propiedad violada por cada enunciado}
\label{tab:violaciones}
\small
\begin{tabular}{C{1.3cm} L{4.0cm} L{8.0cm}}
\toprule
\textbf{Código} & \textbf{Propiedad violada} & \textbf{Motivo} \\
\midrule
E03 & Ambigüedad y no verificabilidad &
La expresión ``rápida'' es subjetiva, porque no establece
un tiempo de respuesta ni una condición de prueba. \\

E07 & No verificabilidad &
La expresión ``seguros'' no especifica medidas de protección
ni criterios que permitan comprobar el cumplimiento. \\

E10 & Ambigüedad y no verificabilidad &
La expresión ``fácil de usar'' depende de la percepción
del usuario y no establece una tarea ni una medida para
evaluar la facilidad de uso. \\

E14 & No verificabilidad e inconsistencia &
La expresión ``siempre disponible'' no establece un porcentaje
ni un periodo de medición. Además, puede generar una
inconsistencia con E19, que admite un plazo de hasta
30 minutos para restablecer el servicio ante una falla. \\

E18 & Ambigüedad y no verificabilidad &
La expresión ``sin problemas'' no establece condiciones
concretas para comprobar que la integración funciona
correctamente. \\
\bottomrule
\end{tabular}
\end{table}

Para mejorar la precisión de estos requisitos se propone
una métrica, una unidad y un umbral que permitan comprobar
su cumplimiento. Los valores planteados son criterios
iniciales de especificación y deben validarse con los
responsables del policlínico antes de considerarse definitivos.
No representan valores establecidos directamente por las
normas ISO o IEEE.

\subsection{Requisitos reformulados}

\begin{enumerate}[leftmargin=1.2cm]

    \item \textbf{E03 -- Eficiencia del desempeño}

    \textbf{Enunciado original:} La aplicación debe ser rápida.

    \textbf{Requisito reformulado:} La aplicación deberá responder
    en un tiempo máximo de 2 segundos para al menos el 95\,\%
    de las operaciones de consulta y reserva, durante una prueba
    con 300 usuarios concurrentes.

    \textbf{Métrica:} Percentil 95 del tiempo de respuesta.

    \textbf{Unidad:} Segundos.

    \textbf{Umbral:} Menor o igual a 2 segundos.

    \item \textbf{E07 -- Seguridad}

    \textbf{Enunciado original:} Los datos de salud de los
    pacientes deben estar seguros.

    \textbf{Requisito reformulado:} El 100\,\% de los accesos
    a los datos de salud deberá requerir autenticación y
    autorización según el rol del usuario. Además, el 100\,\%
    de las comunicaciones que transmitan estos datos deberá
    utilizar un canal cifrado.

    \textbf{Métrica:} Porcentaje de accesos protegidos y de
    comunicaciones cifradas.

    \textbf{Unidad:} Porcentaje.

    \textbf{Umbral:} 100\,\% de cumplimiento en ambos criterios.

    \item \textbf{E10 -- Capacidad de interacción}

    \textbf{Enunciado original:} La plataforma debe ser fácil
    de usar para los pacientes adultos mayores.

    \textbf{Requisito reformulado:} Al menos el 90\,\% de los
    participantes mayores de 60 años que realicen una prueba
    de uso deberá completar el proceso de reserva de una cita
    en un tiempo máximo de 3 minutos, sin recibir asistencia
    durante la ejecución de la tarea.

    \textbf{Métrica:} Porcentaje de participantes que completan
    correctamente la tarea y tiempo de ejecución.

    \textbf{Unidad:} Porcentaje y minutos.

    \textbf{Umbral:} Al menos 90\,\% de éxito y máximo 3 minutos.

    \item \textbf{E14 -- Fiabilidad}

    \textbf{Enunciado original:} La plataforma debe estar
    siempre disponible.

    \textbf{Requisito reformulado:} La plataforma deberá mantener
    una disponibilidad mensual mínima del 99,5\,\%, calculada
    sobre el tiempo total de cada mes. Ante una falla del servidor
    principal, el servicio deberá restablecerse en menos de
    30 minutos, de acuerdo con lo establecido en E19. El tiempo
    durante el cual el servicio permanezca interrumpido se
    contabilizará como indisponibilidad.

    \textbf{Métrica:} Porcentaje mensual de disponibilidad.

    \textbf{Unidad:} Porcentaje.

    \textbf{Umbral:} Igual o superior al 99,5\,\%.

    \item \textbf{E18 -- Compatibilidad}

    \textbf{Enunciado original:} La plataforma debe integrarse
    sin problemas con el sistema de facturación actual.

    \textbf{Requisito reformulado:} La plataforma deberá
    intercambiar correctamente los registros de cobro con
    el sistema de facturación existente, de manera que al
    menos el 99\,\% de las transacciones válidas de una prueba
    de integración se procese correctamente y en un tiempo
    máximo de 5 segundos por transacción.

    \textbf{Métrica:} Porcentaje de transacciones procesadas
    correctamente y tiempo de procesamiento.

    \textbf{Unidad:} Porcentaje y segundos.

    \textbf{Umbral:} Al menos 99\,\% de éxito y máximo 5 segundos.

\end{enumerate}

\subsection{Resumen de métricas y umbrales}

\begin{table}[htbp]
\centering
\caption{Resumen de los requisitos de calidad cuantificados}
\label{tab:metricas}
\small
\begin{tabular}{C{1.3cm} L{3.0cm} L{4.0cm} L{2.0cm} L{2.8cm}}
\toprule
\textbf{Código} &
\textbf{Característica} &
\textbf{Métrica} &
\textbf{Unidad} &
\textbf{Umbral} \\
\midrule

E03 & Eficiencia del desempeño &
Percentil 95 del tiempo de respuesta &
Segundos &
$\leq 2$ \\

E07 & Seguridad &
Accesos y comunicaciones protegidos &
Porcentaje &
100\,\% \\

E10 & Capacidad de interacción &
Éxito de la tarea y tiempo de ejecución &
Porcentaje y minutos &
$\geq 90\,\%$ y $\leq 3$ \\

E14 & Fiabilidad &
Disponibilidad mensual &
Porcentaje &
$\geq 99{,}5\,\%$ \\

E18 & Compatibilidad &
Transacciones correctas y tiempo de procesamiento &
Porcentaje y segundos &
$\geq 99\,\%$ y $\leq 5$ \\

\bottomrule
\end{tabular}
\end{table}

La cuantificación aporta mayor objetividad a los requisitos
y facilita la evaluación de su cumplimiento. Para que una
condición sea verificable, debe contar con criterios que
puedan comprobarse mediante una métrica y un procedimiento
de prueba definido, de acuerdo con ISO/IEC/IEEE 29148:2018
\cite{ieee29148}.

% COMPARACIÓN
\section{Comparación de resultados con un compañero}

En esta sección se registran las diferencias identificadas
al contrastar la clasificación individual con la de un compañero.
Para resolverlas se revisan el contenido de cada caso,
su propósito y los criterios establecidos en las normas
de referencia.

\subsection{Registro de discrepancias}

\begin{longtable}{
    C{0.8cm}
    L{3.0cm}
    L{3.0cm}
    L{3.0cm}
    L{4.0cm}
}
\caption{Registro de discrepancias y resolución}
\label{tab:discrepancias}\\

\toprule
\textbf{N.º} &
\textbf{Código} &
\textbf{Mi clasificación} &
\textbf{Clasificación del compañero} &
\textbf{Resolución y justificación} \\
\midrule
\endfirsthead

\toprule
\textbf{N.º} &
\textbf{Código} &
\textbf{Mi clasificación} &
\textbf{Clasificación del compañero} &
\textbf{Resolución y justificación} \\
\midrule
\endhead

\bottomrule
\endfoot

1 & [Completar] & [Completar] & [Completar] & [Completar] \\

2 & [Completar] & [Completar] & [Completar] & [Completar] \\

3 & [Completar] & [Completar] & [Completar] & [Completar] \\

\end{longtable}

\subsection{Resultado de la comparación}

Se compararon los 20 enunciados con [nombre del compañero]
y se encontraron [n] diferencias. Las discrepancias se
resolvieron revisando el contenido de cada caso y aplicando
los criterios de clasificación establecidos en la guía
FORM-TL-01 y en las normas de referencia.

La clasificación definitiva se presenta en la
Tabla~\ref{tab:clasificacion}. Las diferencias identificadas
y los acuerdos alcanzados se documentan en la
Tabla~\ref{tab:discrepancias}.

% DECLARACIÓN DE USO DE IA

\section{Declaración de uso de herramientas de apoyo}


% REFERENCIAS
\bibliographystyle{ieeetr}
\bibliography{referencias}

\end{document}